\section{Problem Formulation and Methods}\label{sec:method}

\begin{figure*}[!t]
\centering
\maybefigure{architecture}{\textwidth}
\caption{System architecture. The data partition, the DP-SGD step, the aggregation and the privacy accounting are
identical in every run; only the personalisation mechanism changes. The accountant calibrates the noise so that every
client and every mechanism spends the same budget $(\eps_0,\delta)$.}
\label{fig:architecture}
\end{figure*}

\subsection{Federated Setting}
A federation has $K$ clients. Client $k$ holds a private training set $D_k$ of $n_k$ records drawn from its own
distribution, and $n=\sum_k n_k$. Because the client distributions differ, the data are not independent and
identically distributed (non-IID). With $\ell(w;z)$ the loss of parameters $w$ on record $z$ and
$F_k(w)=\frac{1}{n_k}\sum_{z\in D_k}\ell(w;z)$ the local risk, federated averaging (FedAvg)~\cite{mcmahan2017communication}
seeks one global model
\begin{equation}
\min_{w}\; F(w)=\sum_{k=1}^{K} p_k F_k(w), \qquad p_k=\frac{n_k}{n}. \label{eq:fedavg}
\end{equation}
In every round the server sends the current model to the clients, each client runs $S$ local steps, and the server
averages the returned updates with weights $p_k$. All $K$ clients take part in every round.

\subsection{Record-Level Privacy and a Common Budget}
\emph{Threat model.} The server, other clients and anyone who can query a released model (global or personalised)
are treated as potential adversaries; the protected unit is one record of one client. We use record-level rather
than client-level DP for two reasons: each client in the cross-silo setting we model holds the records of many
individuals, and client-level guarantees only become useful with thousands of clients per
round~\cite{kairouz2021advances}, far more than the 20 clients studied here.

\emph{Private gradient query.} Every local step is a DP-SGD step~\cite{abadi2016deep}. Each record of $D_k$ joins the
batch $B$ independently with probability $q$ (Poisson sampling); its gradient $g_i=\nabla\ell(w;z_i)$ is clipped to
norm $C$; Gaussian noise is added to the sum, which is divided by the expected batch size:
\begin{equation}
\tilde g=\frac{1}{q\,n_k}\Big(\sum_{i\in B} g_i\,\min\!\big(1,\tfrac{C}{\lVert g_i\rVert_2}\big)+\mathcal N\!\big(0,\sigma^2C^2 I\big)\Big).
\label{eq:dpsgd}
\end{equation}
Changing one record changes the clipped sum by at most $C$, so the noise level is set by the noise multiplier
$\sigma$ alone.

\emph{Private loss report.} Two of the methods need a client's average training loss: q-FFL to weight clients and
IFCA to choose a cluster. Releasing it is also a query on private data, so it is privatised in the same way: per-record
losses are clipped to $[0,L_{\max}]$ and Gaussian noise with standard deviation $\kappa\sigma L_{\max}$ is added to
their sum before dividing by $n_k$. We use $\kappa=30$, which keeps the noise on the mean loss small for clients with
a few thousand records while charging little of the budget.

\emph{Accounting and calibration.} A client that answers $T_g$ gradient queries and $T_\ell$ loss reports composes a
subsampled Gaussian mechanism with a plain Gaussian mechanism. Its R\'enyi-DP curve~\cite{mironov2017renyi} is
\begin{align}
\rho(\alpha)&=T_g\,\rho_{\mathrm{SG}}(\alpha;q,\sigma)+T_\ell\,\frac{\alpha}{2\kappa^2\sigma^2}, \nonumber\\
\eps&=\min_{\alpha>1}\Big(\rho(\alpha)+\frac{\log(1/\delta)}{\alpha-1}\Big),
\label{eq:accounting}
\end{align}
where $\rho_{\mathrm{SG}}$ is the R\'enyi divergence of one subsampled Gaussian step, computed by the Opacus
accountant~\cite{yousefpour2021opacus}. For every run we find, by bisection, the smallest $\sigma$ whose total
$\eps$ does not exceed the target $\eps_0$. Because all clients take part in every round and run the same number of
steps, every client ends at the same $(\eps_0,\delta)$; because the calibration counts every query a mechanism makes,
every mechanism does too. Personalised models are trained with the same DP-SGD steps, so every released model---global
or personal---carries the guarantee, and evaluation and model selection are post-processing.

\subsection{Personalisation Mechanisms}
Each mechanism produces a model $v_k$ for every client (Fig.~\ref{fig:architecture}). Table~\ref{tab:mechanisms} lists what each one shares and
how it spends its query budget.

\emph{DP-FedAvg (baseline).} Plain FedAvg with DP-SGD local steps; every client uses the global model, $v_k=w$.

\emph{Local-only (baseline).} Each client trains alone for the same number of DP-SGD steps and shares nothing.

\emph{Fine-tuning.} The global model is trained with DP-FedAvg, and each client then adapts it with further DP-SGD
steps on its own data. To keep the step count equal, fine-tuning branches from the FedAvg run after $R-r$ rounds and
spends the remaining $r\cdot S$ steps (10\% of the budget) locally; the global-model and fine-tuning configurations
therefore issue exactly the same number of queries.

\emph{Per-FedAvg (meta-learning).} The first-order variant of~\cite{fallah2020personalized}: each meta-step takes a
DP gradient step with rate $\alpha$ on one batch and evaluates a second DP gradient at the adapted point, which
updates the starting point with rate $\beta$. A meta-step costs two queries, so a client runs $S/2$ meta-steps per
round. After training, each client personalises with a few DP steps of rate $\alpha$, which are also charged.

\emph{Clustered FL (IFCA-style).} The server keeps $m=3$ cluster models~\cite{ghosh2022efficient}. In every round each
client reports a private loss for each cluster model, joins the one with the lowest noisy loss, and trains it; the
server averages updates within each cluster. The client's personalised model is its cluster model.

\emph{FedPer (parameter decoupling).} The network is split into a shared body and a client-specific
head~\cite{arivazhagan2019federated}; only the body is averaged, and the head never leaves the client. The head is
trained with the same DP-SGD steps, so it is protected like every other parameter.

\begin{table}[!t]
\caption{Mechanisms under one budget: what is exchanged and what is charged per client and round ($P$ = model
parameters, $P_b$ = body parameters, $S$ = local steps, $m$ = clusters).}
\label{tab:mechanisms}
\centering
\footnotesize
\begin{tabular}{p{1.45cm}p{1.75cm}p{1.6cm}p{2.2cm}}
\toprule
Mechanism & Down / up per client & Private queries & Personal model \\
\midrule
DP-FedAvg & $P$ / $P$ & $S$ & global model \\
Local-only & 0 / 0 & $S$ & own model \\
Fine-tuning & $P$ / $P$ & $S$ (last 10\% local) & global + local steps \\
Per-FedAvg & $P$ / $P$ & $S$ (+ adaptation) & meta-model + adaptation \\
IFCA & $mP$ / $P$ & $S$ + $m$ loss reports & chosen cluster model \\
FedPer & $P_b$ / $P_b$ & $S$ & shared body + own head \\
q-FFL & $P$ / $P{+}1$ & $S$ + 1 loss report & global model \\
\bottomrule
\end{tabular}
\end{table}

\subsection{Fairness-Constrained Aggregation (q-FFL)}
q-FFL~\cite{li2020fair} replaces (\ref{eq:fedavg}) with $\min_w\sum_k \frac{p_k}{q+1}F_k^{q+1}(w)$; a larger $q$ gives
more weight to clients with high loss and $q=0$ recovers FedAvg. We use the q-FedAvg update of~\cite{li2020fair}: with
local model $\bar w_k$ after $S$ DP steps, $\Delta_k=L(w-\bar w_k)$ for $L=1/\eta$ and the private loss report
$\tilde F_k$, the server sets
\begin{equation}
w \leftarrow w-\frac{\sum_k \tilde F_k^{\,q}\Delta_k}{\sum_k\big(q\tilde F_k^{\,q-1}\lVert\Delta_k\rVert^2+L\tilde F_k^{\,q}\big)}.
\label{eq:qffl}
\end{equation}
The loss reports are the only extra queries, and they are charged through (\ref{eq:accounting}); at a fixed $\eps_0$
the calibration therefore adds slightly more gradient noise, which is how any ``privacy price'' of q-FFL shows up.

\subsection{Evaluation Metrics}
Let $a_k$ be the test accuracy of client $k$'s model on its own held-out data, and $A_k(\cdot)$ the accuracy of a
given model on that data.

\emph{Utility and personalisation gain.} We report the mean accuracy $\bar a$ and the gain of each mechanism $M$ over
DP-FedAvg trained under the same budget,
\begin{equation}
G_k(M,\eps)=A_k\big(v_k(M,\eps)\big)-A_k\big(w(\eps)\big). \label{eq:gain}
\end{equation}
The share of clients with $G_k<0$ counts the clients that personalisation makes worse off.

\emph{Fairness across clients.} The worst-10\% accuracy is the mean of the $\lceil 0.1K\rceil$ lowest $a_k$; we also
report the standard deviation of $a_k$ and Jain's index~\cite{jain1984quantitative},
\begin{equation}
J=\frac{\big(\sum_{k}a_k\big)^2}{K\sum_{k}a_k^2}, \label{eq:jain}
\end{equation}
which is 1 when all clients have equal accuracy and approaches $1/K$ when one client receives everything.

\emph{Group fairness} (Adult only). With $\hat y$ the prediction and $s$ the sensitive attribute, the
demographic-parity gap is $\lvert P(\hat y{=}1\mid s{=}1)-P(\hat y{=}1\mid s{=}0)\rvert$ and the equal-opportunity gap
is the same difference for true-positive rates.

\emph{Cost of privacy.} Accuracy loss per unit budget between neighbouring budgets $\eps_a>\eps_b$ is
$(\bar a(\eps_a)-\bar a(\eps_b))/(\eps_a-\eps_b)$, reported with the drop relative to the non-private run.

\emph{Measured leakage.} $\eps$ is a worst-case bound, not a measurement. We therefore mount a black-box
loss-threshold membership-inference attack~\cite{yeom2018privacy} on every client's model, with a random subset of
its training records as members and an equal number of its held-out records as non-members, and report the area
under the ROC curve (AUC; 0.5 is random guessing) and the attack advantage (true-positive minus false-positive rate)
at the threshold equal to the mean training loss.

\emph{System cost.} Bytes exchanged per round and in total (32-bit parameters), wall-clock seconds per round, peak
GPU memory, and the time overhead of DP relative to the non-private run.

\subsection{Problem Statement}
For a fixed budget $\eps_0$ and a protocol in which the partition, model, clipping norm, sampling rate, number of
rounds and accounting are identical across mechanisms, we seek the mechanism with the largest mean personalisation
gain that does not make the accuracy distribution less even than DP-FedAvg:
\begin{equation}
M^\star=\arg\max_{M}\;\frac{1}{K}\sum_{k=1}^{K}G_k(M,\eps_0)\quad\text{s.t.}\quad J(M,\eps_0)\ge J_0,
\label{eq:problem}
\end{equation}
where $J_0$ is Jain's index of DP-FedAvg at $\eps_0$. Solving (\ref{eq:problem}) for every budget in the sweep gives
the recommended operating configuration. Because (\ref{eq:problem}) assumes that an equal $\eps$ means comparable
leakage, we test that assumption with the membership-inference attack instead of taking it for granted.
